%%%%%%%%%%%%%%%%%%%%%%%%%%%%%%%%%%%%%%%%%%%%%%
%% preamble.tex --- 版式层 (LAYOUT LAYER)
%%
%% 这是唯一决定"长什么样"的文件：class、边距、段落、配色、
%% 定理环境的外观、表格的线。改这一个文件，所有章节一起变；
%% 各章正文 .tex 里一个版式指令都没有。
%%
%% 骨架来自 Harry 的作业模板；作业专用的 \problem / \answer
%% 计数器换成了笔记要用的编号定理环境。
%%
%% 正文层依赖的接口（替换本文件时必须继续提供）：
%%   环境 definition / claim / example  编号，共用一个计数器
%%   环境 note                          不编号
%%   环境 proof                         amsthm，以 ∎ 收尾
%%   命令 \fig{stem}{caption}{label}    按主题选 -light / -dark 图
%%   命令 \chapref{显示文字}{stem}       指向同目录另一章的 PDF
%%   命令 \secref{label}                § 号交叉引用
%%   环境 checklist + \checkitem        末尾的复选框清单
%%   命令 \chapterhead{号}{标题}{stem}   章首：单章模式出标题页+目录，
%%                                        合订模式出章标题并重编号
%%   开关 \ifdark                       由 build.sh 通过 \DARKMODE 置位
%%   开关 \ifbook                       由 build.sh 通过 \BOOKMODE 置位
%%                                        (all_chapters.tex 合订本)
%%%%%%%%%%%%%%%%%%%%%%%%%%%%%%%%%%%%%%%%%%%%%%

\documentclass[12pt]{article}


%%%%%%%%%%%%%%%%%%%%%%%%%%%%%%%%%%%%%%%%%%%%
%% Page formatting.
%%%%%%%%%%%%%%%%%%%%%%%%%%%%%%%%%%%%%%%%%%%%%%
\usepackage[margin=1in]{geometry} % adjust margins
\setlength\parindent{0em} % never indent a paragraph

% Paragraphs are separated by space rather than by an indent, so the skip has
% to be visible -- but a rigid 1em was a whole blank line between every two
% paragraphs, and a chapter is mostly paragraphs. 0.65em still reads as a
% break; the rubber lets a page that is nearly full close up instead of
% pushing a line over.
\setlength{\parskip}{0.65em plus 0.15em minus 0.15em}

% Space is added at the bottom of a short page rather than stretched between
% its paragraphs -- otherwise \parskip's "plus" opens gaps that look
% deliberate and are not.
\raggedbottom


%%%%%%%%%%%%%%%%%%%%%%%%%%%%%%%%%%%%%%%%%%%%
%% Packages
%%%%%%%%%%%%%%%%%%%%%%%%%%%%%%%%%%%%%%%%%%%%%%
\usepackage{amsthm, amssymb, amsmath} % basic math stuff -- always include this
\usepackage{enumerate}   % customize enumeration
\usepackage[table]{xcolor} % named colours (+ \arrayrulecolor), before hyperref
\usepackage{tcolorbox}   % allows me to make colored boxes
\usepackage{graphicx}    % the chapter figures
\usepackage{tabularx}    % tables whose text column fills the line
\usepackage{xltabular}   % ...and that may break across pages
\usepackage{array}       % >{...} column prefixes
\usepackage{caption}
\usepackage{titlesec}    % heading spacing; article's defaults double up with \parskip
\usepackage{etoolbox}    % \AtBeginEnvironment, to set table type size in one place


%%%%%%%%%%%%%%%%%%%%%%%%%%%%%%%%%%%%%%%%%%%%%%
%% Theme. build.sh defines \DARKMODE for the dark build and
%% nothing for the light one. Page colour, link colour and the
%% table rules all have to agree, or a dark PNG lands in a white
%% margin (and vice versa).
%%%%%%%%%%%%%%%%%%%%%%%%%%%%%%%%%%%%%%%%%%%%%%
\newif\ifdark
\ifdefined\DARKMODE \darktrue \fi

% build.sh defines \BOOKMODE only for all_chapters.tex, where every chapter
% body is \input into one document instead of compiling on its own.
\newif\ifbook
\ifdefined\BOOKMODE \booktrue \fi

\ifdark
  \definecolor{pagebg}{HTML}{14161A}
  \definecolor{ink}{HTML}{E6E8EB}
  \definecolor{linkcol}{HTML}{6FB7E8}
  \definecolor{rulecol}{HTML}{4A5057}
  \definecolor{boxbg}{HTML}{1E2126}
\else
  \definecolor{pagebg}{HTML}{FFFFFF}
  \definecolor{ink}{HTML}{000000}
  \definecolor{linkcol}{HTML}{0000FF}
  \definecolor{rulecol}{HTML}{000000}
  \definecolor{boxbg}{HTML}{F2F4F7}
\fi

\usepackage[colorlinks=true, allcolors=linkcol]{hyperref} % clickable hyperlinks

\ifdark
  \usepackage{pagecolor}
  \pagecolor{pagebg}
  \color{ink}
  \arrayrulecolor{rulecol}
\fi


%%%%%%%%%%%%%%%%%%%%%%%%%%%%%%%%%%%%%%%%%%%%%%
%% Type scale and vertical rhythm.
%%
%% article's section skips (3.5ex before, 2.3ex after) were written for a
%% document whose paragraphs are marked by an indent and carry no \parskip.
%% Here every paragraph already brings 0.65em with it, so the defaults land
%% on top of that and a heading floats in the middle of its own white band.
%% These are the totals we actually want.
%%%%%%%%%%%%%%%%%%%%%%%%%%%%%%%%%%%%%%%%%%%%%%
\titlespacing*{\section}      {0pt}{1.1em plus 0.2em minus 0.2em}{0.35em}
\titlespacing*{\subsection}   {0pt}{0.9em plus 0.2em minus 0.2em}{0.25em}
\titlespacing*{\subsubsection}{0pt}{0.7em plus 0.2em minus 0.2em}{0.2em}

% Tables carry the reference material -- long cells, many rows. One step down
% fits more of a comparison on the page without touching the prose size, and
% the caption goes with it so a figure's label never outweighs its own axis
% labels.
\AtBeginEnvironment{tabularx}{\small}
\AtBeginEnvironment{xltabular}{\small}
\AtBeginEnvironment{tabular}{\small}
\captionsetup{font=small, labelfont=bf, skip=0.5em}

% Floats: article will bail a figure out to a page of its own rather than let
% it take more than 70 percent of a text page. At \linewidth that happens
% constantly, and the bailed-out page is mostly white.
\renewcommand{\topfraction}{0.92}
\renewcommand{\bottomfraction}{0.7}
\renewcommand{\textfraction}{0.07}
\renewcommand{\floatpagefraction}{0.75}
\setlength{\intextsep}{0.9em plus 0.2em minus 0.2em}
\setlength{\textfloatsep}{1.0em plus 0.2em minus 0.2em}


%%%%%%%%%%%%%%%%%%%%%%%%%%%%%%%%%%%%%%%%%%%%%%
%% Numbered theorem environments.
%%
%% The head is run-in and bold, so a Definition still opens the
%% paragraph the way "\textbf{Problem 1.}" does in the homework
%% template --- it just carries a number now, and can be \ref'd.
%% definition / claim / example share one counter, so the numbers
%% read in document order and never collide with the § numbers.
%%%%%%%%%%%%%%%%%%%%%%%%%%%%%%%%%%%%%%%%%%%%%%
\newtheoremstyle{notes}
  {\topsep}    % space above
  {\topsep}    % space below
  {}           % body font
  {0pt}        % indent
  {\bfseries}  % head font
  {.}          % punctuation after head
  { }          % space after head -- a plain space makes it run-in
  {}           % head spec -- empty means the default layout
\theoremstyle{notes}
\newtheorem{definition}{Definition}
\newtheorem{claim}[definition]{Claim}
\newtheorem{example}[definition]{Example}

% Note is a caveat, not a result: unnumbered, nothing refers back to it.
\newtheorem*{note}{Note}

% amsthm sets the proof head in italic; this template wants it bold.
\renewcommand{\proofname}{\normalfont\bfseries Proof}


%%%%%%%%%%%%%%%%%%%%%%%%%%%%%%%%%%%%%%%%%%%%
%% Define (or redefine) any commands/macros here
%%%%%%%%%%%%%%%%%%%%%%%%%%%%%%%%%%%%%%%%%%%%%%
\newcommand{\R}{\mathbb{R}} % sets up the command \R to write the R symbol for real numbers
\newcommand{\Z}{\mathbb{Z}} % use \Z for integers
\newcommand{\Q}{\mathbb{Q}} % use \Q for rationals
\newcommand{\N}{\mathbb{N}} % use \N for natural numbers

% --- notes-specific -------------------------------------------------

% \fig{stem}{caption}{label} -- picks <name>-light.png or -dark.png to match
% the build, so the figure and the page always agree.
%
% Every chapter's figures live beside their generator in figures/<NN>_figures/.
% LaTeX is run from docs/, so these are the paths it searches, in order; PNG
% names are unique across chapters, so a chapter never has to say which
% directory its own figure is in.
\graphicspath{%
  {figures/00_figures/}{figures/01_figures/}{figures/02_figures/}%
  {figures/03_figures/}{figures/04_figures/}{figures/05_figures/}%
  {figures/09_figures/}}
\newcommand{\figvariant}{light}
\ifdark \renewcommand{\figvariant}{dark} \fi
\newcommand{\fig}[3]{%
  \begin{figure}[htbp]
    \centering
    \includegraphics[width=\linewidth]{#1-\figvariant}
    \caption{#2}\label{#3}
  \end{figure}}

% \chapterhead{number}{title}{stem} -- opens a chapter.
%
% Standalone: a title page and the chapter's own table of contents, which is
% what a single-chapter PDF in read_only_chapters/ wants.
%
% Book: a chapter heading inside the running document.
%
% Either way the chapter number is carried into every counter below it:
% sections print as "05.1", subsections as "05.1.2", and the shared theorem
% counter restarts too, so Definition 3 of chapter 05 reads "Definition 05.3"
% and cannot collide with chapter 02's. A section number therefore means the
% same thing in the standalone PDF and in the bound edition, and \secref
% prints "§ 05.4.3" in both.
\newcommand{\chapterhead}[3]{%
  \ifbook
    \clearpage
    \phantomsection\label{chap:#3}%
    \addcontentsline{toc}{part}{#1 \midot{} #2}%
    {\huge\bfseries #1 \midot{} #2\par}\vspace{0.4em}%
  \else
    \title{#1 \midot{} #2}%
    \author{Harry Chen}%
    \date{}%
    \maketitle
    \tableofcontents
    \newpage
  \fi
  \setcounter{section}{0}%
  \renewcommand{\thesection}{#1.\arabic{section}}%
  \setcounter{definition}{0}%
  \renewcommand{\thedefinition}{#1.\arabic{definition}}%
  \setcounter{figure}{0}%
  \renewcommand{\thefigure}{#1.\arabic{figure}}%
}

% \chapref{text}{stem} -- a cross-chapter reference.
%
% Standalone: an href to the sibling PDF, which sits beside this one in
% read_only_chapters/. Book: a jump to that chapter's \chapterhead, or plain
% text when the chapter is not in the book yet (still Markdown, say) -- an
% unconverted chapter must not print "??".
\newcommand{\chapref}[2]{%
  \ifbook
    \ifcsname r@chap:#2\endcsname\hyperref[chap:#2]{#1}\else#1\fi
  \else
    \href{#2.pdf}{#1}%
  \fi}

% \secref{label} -- "§ 4.3", clickable. This replaces the hand-typed
% section numbers the Markdown had to carry.
\newcommand{\secref}[1]{\S\ref{#1}}

% A forward pointer to another chapter or to code.
\newcommand{\pointer}[1]{\par\noindent$\rightarrow$~#1\par}

% The closing checklist.
\newenvironment{checklist}{\begin{itemize}\setlength{\itemsep}{0pt}}{\end{itemize}}
\newcommand{\checkitem}{\item[$\square$]}

% Interpunct, for titles like "05 · Modelling".
\newcommand{\midot}{\textperiodcentered}

%%%%%%%%%%%%%%%%%%%%%%%%%%%%%%%%%%%%%%%%%%%%%%

%%%%%%%%%%%%%%%%%%%%%%%%%%%%%%%%%%%%%%%%%%%%%%
%% Two consequences of \parskip, fixed.
%%%%%%%%%%%%%%%%%%%%%%%%%%%%%%%%%%%%%%%%%%%%%%

% The table of contents is a list, not a run of paragraphs; at 1em of
% skip per entry it spills onto a second page. Zero the skip for its
% duration only, so the body keeps the spacing the template wants.
\makeatletter
\let\notes@toc\tableofcontents
\renewcommand{\tableofcontents}{%
  \begingroup\setlength{\parskip}{0pt}\notes@toc\endgroup}
\makeatother

% \parskip does not reach inside a tabular, so rows come out tighter
% than the surrounding text. Open them back up.
\renewcommand{\arraystretch}{1.25}

% article sizes the table-of-contents number columns for "1" and "1.2".
% \chapterhead prints "05.1" and "05.4.3", which overrun them and collide
% with the title text, so each column is widened and the entries re-indented.
\makeatletter
\renewcommand*\l@section{\@dottedtocline{1}{1.5em}{3.2em}}
\renewcommand*\l@subsection{\@dottedtocline{2}{4.7em}{4.0em}}
\renewcommand*\l@subsubsection{\@dottedtocline{3}{8.7em}{4.8em}}
\makeatother
%%%%%%%%%%%%%%%%%%%%%%%%%%%%%%%%%%%%%%%%%%%%%%
